\documentclass[11pt,a4paper]{article}

\usepackage[utf8]{inputenc}
\usepackage[T1]{fontenc}
\usepackage{textcomp}
\usepackage[french]{babel}
\addto\captionsfrench{\renewcommand{\tablename}{Tableau}\renewcommand{\figurename}{Figure}}

\usepackage{lmodern}\usepackage{geometry}\usepackage{setspace}
\usepackage{microtype}\usepackage{xcolor}
\usepackage{booktabs}\usepackage{tabularx}\usepackage{array}\usepackage{enumitem}
\usepackage{titlesec}\usepackage{fancyhdr}\usepackage{caption}\usepackage{float}
\usepackage{amssymb}\usepackage{amsmath}\usepackage[most]{tcolorbox}\usepackage{tikz}
\usetikzlibrary{positioning,arrows.meta,calc}
\usepackage{adjustbox}
\usepackage{hyperref}
\geometry{top=2.1cm,bottom=2.1cm,left=2.2cm,right=2.1cm}
\setstretch{1.05}
\setlength{\emergencystretch}{1.2em}
\setlength{\parskip}{0.3em}

\definecolor{mainblue}{RGB}{28,74,102}
\definecolor{darkgray}{RGB}{62,62,66}
\definecolor{lightblue}{RGB}{238,245,250}
\definecolor{ent}{RGB}{221,236,247}
\definecolor{tab}{RGB}{215,237,221}
\definecolor{mod}{RGB}{233,233,238}
\definecolor{ext}{RGB}{250,233,207}
\definecolor{alerte}{RGB}{250,226,226}
\definecolor{codebg}{RGB}{246,246,248}
\definecolor{kored}{RGB}{176,42,42}

\hypersetup{colorlinks=true,linkcolor=mainblue,urlcolor=mainblue,
  pdftitle={Guide de la solution - workflow, code et technologies},pdfauthor={Equipe Yahya, Sohayb, Younes}}

\titleformat{\section}{\Large\bfseries\color{mainblue}}{\thesection}{0.6em}{}
\titleformat{\subsection}{\large\bfseries\color{darkgray}}{\thesubsection}{0.6em}{}
\titlespacing*{\section}{0pt}{1.2em}{0.45em}
\titlespacing*{\subsection}{0pt}{0.9em}{0.3em}
\captionsetup{font=small,labelfont=bf}
\setlist[itemize]{leftmargin=1.3em,itemsep=0.1em,topsep=0.2em}
\setlist[enumerate]{leftmargin=1.6em,itemsep=0.1em,topsep=0.2em}
\newcolumntype{Y}{>{\raggedright\arraybackslash}X}
\newcolumntype{L}[1]{>{\raggedright\arraybackslash}p{#1}}

\pagestyle{fancy}\fancyhf{}
\lhead{\small\color{darkgray}Guide de la solution --- workflow, code, technologies}
\rhead{\small\color{darkgray}Sujet 1}
\cfoot{\small\thepage}
\renewcommand{\headrulewidth}{0.3pt}

\newtcolorbox{cle}{breakable,colback=lightblue,colframe=mainblue,boxrule=0.7pt,arc=1.4mm,
  left=3mm,right=3mm,top=1.5mm,bottom=1.5mm}
\newtcolorbox{code}{breakable,colback=codebg,colframe=darkgray!35,boxrule=0.4pt,arc=1mm,
  left=2mm,right=2mm,top=0.8mm,bottom=0.8mm}

\DeclareRobustCommand{\f}[1]{{\texttt{\small\renewcommand{\_}{\char`\_\allowbreak}#1}}}
\newcommand{\eur}{\,\texteuro}

\tikzset{
  etage/.style={rectangle,draw=mainblue,fill=ent,rounded corners=2pt,align=center,
                font=\small,inner sep=5pt,minimum height=0.9cm,text width=7.2cm},
  cote/.style={rectangle,draw=darkgray!70,fill=mod,rounded corners=2pt,align=left,
               font=\scriptsize,inner sep=4pt,text width=6.6cm},
  acteur/.style={rectangle,draw=darkgray,fill=ext,rounded corners=2pt,align=center,
                 font=\small,inner sep=5pt,minimum height=0.8cm,text width=7.2cm},
  fl/.style={-{Latex[length=2mm]},semithick,draw=mainblue!80},
  lien/.style={densely dashed,draw=darkgray!60},
}

\begin{document}

\begin{center}
{\small\scshape\color{darkgray} Hackathon IA --- La Banque Postale $\times$ CGI $\times$ ENSEIRB-MATMECA --- Sujet 1\par}
\vspace{0.25cm}
{\LARGE\bfseries\color{mainblue} Guide de la solution\par}
\vspace{0.15cm}
{\large Ce qu'on fait, comment, et dans quel code\par}
\vspace{0.2cm}
{\small Yahya (IA) --- Sohayb (génie logiciel) --- Younes (calcul intensif) --- mardi 27 octobre 2026\par}
{\footnotesize\color{darkgray} État du dépôt : 171 tests et 10 scénarios au vert.\par}
\end{center}

% ════════════════════════════════════════════════════════════════
\section{La solution}

\subsection{En une phrase}

Un client écrit à sa banque ; notre assistant \textbf{comprend} la demande, \textbf{pose les
questions qui manquent}, \textbf{décide} avec des règles bancaires publiques, \textbf{agit} si
besoin (opposition, virement), puis \textbf{répond} en français --- en n'utilisant un LLM qu'au
dernier moment, et seulement si c'est utile.

\subsection{Le principe : une cascade}

Chaque étage ne passe la main au suivant que s'il ne sait pas faire lui-même. Le moins cher
d'abord, le LLM en dernier. C'est l'image donnée par l'organisateur : le vélo pour aller chercher
le pain, la voiture pour le week-end.

\begin{figure}[H]
\centering
\begin{adjustbox}{max width=\textwidth}
\begin{tikzpicture}[node distance=0.45cm]
  \node[acteur] (msg) {Message du client};
  \node[etage,below=of msg] (e1) {\textbf{1. Comprendre} --- de quoi parle-t-il ?};
  \node[etage,below=of e1] (e2) {\textbf{2. Converser} --- que me manque-t-il ?};
  \node[etage,below=of e2] (e3) {\textbf{3. Conseiller} --- quel produit ?\\ pourquoi, pourquoi pas ?};
  \node[etage,below=of e3] (e4) {\textbf{4. Agir} --- opposition, virement, rappel};
  \node[etage,below=of e4] (e5) {\textbf{5. Formuler} --- la réponse en français};
  \node[acteur,below=of e5] (rep) {Réponse\\ + panneaux \og Pourquoi \fg{} et \og Frugalité \fg};
  \node[cote,right=0.7cm of e1] (c1) {réponse attendue ? $\rightarrow$ mots-clés $\rightarrow$ classifieur\\
     rien de sûr : question au client, puis conseiller proposé (et journal)};
  \node[cote,right=0.7cm of e2] (c2) {parcours YAML + profil du client ;\\ on ne redemande jamais ce qu'on sait};
  \node[cote,right=0.7cm of e3] (c3) {règles publiques (35\,\%, 25 ans, plafonds) ;\\ raisons et produits écartés};
  \node[cote,right=0.7cm of e4] (c4) {fonctions sur un SI simulé ;\\ rien sans un \og oui \fg{} du client};
  \node[cote,right=0.7cm of e5] (c5) {gabarit, ou petit LLM une seule fois,\\ sous garde-fou ; mention légale du crédit};
  \foreach \a/\b in {msg/e1,e1/e2,e2/e3,e3/e4,e4/e5,e5/rep} \draw[fl] (\a) -- (\b);
  \foreach \i in {1,...,5} \draw[lien] (e\i) -- (c\i);
\end{tikzpicture}
\end{adjustbox}
\caption{La cascade. Selon la demande, l'étage 3 ou l'étage 4 est sauté.}
\end{figure}

\subsection{Un exemple complet : Léa veut un appartement}

\begin{center}\small
\begin{tabularx}{\textwidth}{@{}L{3.6cm} Y@{}}
\toprule
\textbf{Message} & \textbf{Ce qui se passe} \\
\midrule
\og Je veux acheter un appartement \fg &
  \textbf{Comprendre} : les mots-clés \og acheter un appartement \fg{} donnent
  \f{credit\_immobilier}, en moins d'une milliseconde, sans modèle.
  \textbf{Converser} : le parcours \f{credit\_immobilier.yaml} demande 4 informations ; revenus et
  crédits en cours sont \emph{dans son profil}, donc confirmés ; il manque le prix et l'apport.
  \textbf{Formuler} : \og C'est un beau projet. J'ai noté vos revenus de 2\,900\eur{} nets par
  mois. Quel est le prix du bien\dots{} ? \fg{} (gabarit, sans LLM). \\
\addlinespace
\og Autour de 250\,000, et je peux mettre les 25\,000 \fg &
  \textbf{Comprendre} : une question était en attente, donc c'est \emph{la réponse} ; les
  montants sont extraits par règles. \textbf{Conseiller} : règles HCSF au taux moyen Banque de France de 3,32\,\% : mensualité de
  1\,197\eur{}, soit \textbf{41,3\,\% des revenus, au-delà des 35\,\%} : projet non retenu, sans refus
  bancaire ; crédit sur 30 ans écarté (25 ans maximum) ; budget indicatif d'environ 215\,000\eur{} ;
  hypothèses et sources affichées ; rendez-vous conseiller proposé.
  \textbf{Formuler} : explication, mention légale, proposition de rappel. \\
\addlinespace
\og Et avec un bien à 200\,000 ? \fg &
  Compris comme une \textbf{correction} : le prix est mis à jour et le conseil recalculé :
  32,2\,\%, orientation vers un conseiller. \\
\bottomrule
\end{tabularx}
\end{center}

Pendant ce temps, l'interface affiche l'intention et comment elle a été comprise, les règles
appliquées, les produits proposés et écartés, et le temps de chaque étape --- sans aucun appel au
LLM.

\subsection{Chaque étage : quoi, comment, pourquoi}

\begin{center}\small
\begin{tabularx}{\textwidth}{@{}L{2.4cm} Y L{3.6cm}@{}}
\toprule
\textbf{Étage} & \textbf{Comment} & \textbf{Pourquoi} \\
\midrule
1. Comprendre \newline (Younes) &
  (1) réponse à la question en cours ; (2) mots-clés, qui ne concluent que s'ils sont sûrs ;
  (3) classifieur \f{multilingual-e5-small} + régression logistique, entraîné sur Banking77
  \textbf{en anglais, sans traduction} (à faire) ; (4) sinon : question au client, puis conseiller
  proposé et message noté dans un journal. Pas de gros modèle par défaut. &
  Comprendre ne demande pas de LLM : rapide, presque gratuit, mesurable. \\
\addlinespace
2. Converser \newline (Yahya) &
  Un parcours par situation, en YAML (quoi demander, où le trouver, quoi faire à la fin) ; un
  automate les lit. Cas limites : \og je ne sais pas \fg, correction, abandon. &
  Déterministe, lisible, testable, sans LLM. \\
\addlinespace
3. Conseiller \newline (Sohayb) &
  Règles publiques sourcées : taux d'effort $\leq$ 35\,\%, durée $\leq$ 25 ans, plafonds des
  livrets, âges. Chaque proposition a ses raisons, chaque produit écarté la sienne. &
  Une décision financière doit être juste, reproductible et explicable. \\
\addlinespace
4. Agir \newline (Sohayb) &
  5 actions sur un SI simulé : opposition, virement, explication des frais, rappel conseiller,
  lecture du profil. Tout effet de bord exige un \og oui \fg. &
  Sécurité du client. \\
\addlinespace
5. Formuler \newline (Yahya) &
  Gabarits par défaut. Si activé, le petit LLM (\f{ministral-3:3b}) reformule un conseil, une
  seule fois ; le garde-fou refuse tout nombre absent du gabarit. Mode \og mots simples \fg{} pour
  les clients âgés ou peu à l'aise. &
  Exact et gratuit d'abord ; le LLM seulement pour le naturel. \\
\addlinespace
Mesurer \newline (Younes) &
  Chaque étage est chronométré, les appels au LLM comptés ; CodeCarbon à venir pour comparer
  A (tout LLM), B (notre cascade) et C (sans LLM). &
  La frugalité vaut 20\,\% de la note. \\
\bottomrule
\end{tabularx}
\end{center}

\subsection{Les données}

\begin{center}\small
\begin{tabularx}{\textwidth}{@{}L{4.6cm} Y@{}}
\toprule
\textbf{Donnée} & \textbf{Rôle} \\
\midrule
Banking77 (anglais, 77 intentions) & Entraîner le classifieur, \textbf{sans traduction} \\
Bank Marketing (45\,211 clients) & Ancrer les 8 clients de démo ; futur modèle d'appétence \\
\f{clients.json} & Inès, Marc, Sophie, Paul, Karim, Julie, Léa, Lucas \\
\f{catalogue.json} & Produits et paramètres réglementaires, avec leur source \\
\f{glossaire.json} & Explications du jargon bancaire \\
\bottomrule
\end{tabularx}
\end{center}

\subsection{Où on en est}

\textbf{Fait} : toute la chaîne, les 9 parcours, les règles, les actions, l'interface, le mode
\og mots simples \fg, le journal des incompris, le branchement réel du LLM, la démo de secours et
le support de soutenance --- 171 tests et 10 scénarios au vert.
\textbf{À faire} : classifieur (Younes) ; choix du LLM sur un poste de la salle (Yahya) ;
appétence et sources du catalogue (Sohayb) ; mesures d'énergie, vidéo, diapositives.

\begin{cle}
\textbf{En 30 secondes.} \emph{Notre assistant comprend avec de simples mots-clés ou un petit
modèle de 118 millions de paramètres, décide avec des règles bancaires que le client peut lire,
et n'appelle un LLM qu'une fois, à la fin, pour mettre en mots ce qui est déjà calculé. Chaque
conseil dit pourquoi, et pourquoi pas. Et on l'a mesuré face à une version où un gros LLM fait
tout.}
\end{cle}

% ════════════════════════════════════════════════════════════════
\section{Le code}

\subsection{Les technologies}

\begin{center}\small
\begin{tabularx}{\textwidth}{@{}L{3.6cm} L{4.2cm} Y@{}}
\toprule
\textbf{Besoin} & \textbf{Technologie} & \textbf{Remarque} \\
\midrule
Langage & Python 3.10+ & Tout le projet \\
Interface & Streamlit & Une page web écrite en Python \\
Parcours de conversation & YAML (PyYAML) & Fichiers lisibles par un humain \\
Données & JSON & Clients, catalogue, glossaire \\
LLM local & Ollama (poste de la salle) ou vLLM (DeepEirb, API compatible OpenAI) & Appelés avec \f{urllib} (bibliothèque standard) ; serveur choisi par rôle ; désactivé par défaut \\
Tests & pytest + \f{streamlit.testing} & L'interface est testée sans navigateur \\
À venir (Younes) & sentence-transformers, scikit-learn, CodeCarbon & Classifieur et mesure d'énergie \\
\bottomrule
\end{tabularx}
\end{center}

Dépendances installées aujourd'hui : \f{streamlit}, \f{pyyaml}, \f{pytest}. Tout le reste vient
de la bibliothèque standard de Python.

\subsection{L'arborescence}

\begin{code}
\footnotesize
\begin{verbatim}
hackathon-IA/
|-- app.py                  Interface Streamlit                                   (Yahya)
|-- src/
|   |-- contrats.py         La langue commune : structures echangees entre modules
|   |-- taxonomie.py        14 familles, 85 intentions (77 Banking77 + 8 a nous)
|   |-- assistant.py        reply() : l'orchestrateur, un tour de conversation    (Yahya)
|   |-- llm.py              Appel a Ollama (desactive par defaut)                 (Yahya)
|   |-- texte.py            Normaliser un texte, ecrire 1 197 EUR et 41,3 %
|   |-- incompris.py        Journal des messages non compris                      (Yahya)
|   |-- comprendre/         ETAGE 1                                               (Younes)
|   |   |-- __init__.py     comprendre() : reponse attendue -> mots-cles -> classifieur
|   |   |-- mots_cles.py    Lexique par intention + regle "sur ou pas"
|   |   |-- entites.py      Regex : montants, durees, liens, oui/non
|   |   |-- classifieur.py  BOUCHON : e5-small + regression logistique, a entrainer
|   |   `-- escalade.py     Appel au 14B en dernier recours (desactive par defaut)
|   |-- converser/          ETAGES 2 et 5                                         (Yahya)
|   |   |-- dialogue.py     L'automate : lit les parcours, decide quoi faire
|   |   |-- parcours/*.yaml 9 parcours, un par situation du sujet
|   |   |-- formuler.py     Gabarits, LLM eventuel, mode "mots simples"
|   |   |-- garde_fou.py    Refuse les nombres inventes et les promesses
|   |   `-- glossaire.py    Explication d'un terme bancaire
|   |-- conseiller/         ETAGES 3 et 4                                         (Sohayb)
|   |   |-- clients.py      Charge les profils (adaptateur de la base clients)
|   |   |-- catalogue.py    Charge produits et parametres sources
|   |   |-- regles.py       Moteur de regles : credit, epargne, pret etudiant
|   |   `-- outils.py       Les 5 actions + confirmation obligatoire + journal
|   `-- mesure/__init__.py  Chronometrage de chaque etape                        (Younes)
|-- data/raw/               Banking77, Bank Marketing (bruts)
|-- data/processed/         clients.json, catalogue.json, glossaire.json
|-- eval/scenarios/*.yaml   9 conversations de reference
|-- eval/run_scenarios.py   Les rejoue et verifie les attentes
|-- eval/bench_formulation.py   Choisit le LLM en descendant en taille
`-- tests/                  171 tests pytest
\end{verbatim}
\end{code}

\subsection{Le contrat : comment les modules se parlent}

Les modules ne s'appellent pas n'importe comment : ils échangent des \textbf{dataclasses}
définies une seule fois dans \f{src/contrats.py}. Chacun peut ainsi avancer sans attendre les
autres.

\begin{center}\small
\begin{tabularx}{\textwidth}{@{}L{3.3cm} L{2.6cm} Y@{}}
\toprule
\textbf{Structure} & \textbf{Rendue par} & \textbf{Contenu} \\
\midrule
\f{Comprehension} & Comprendre & famille, intention, confiance, niveau (mots-clés, classifieur,
  réponse attendue, escalade), entités (montants, durée, lien, objet), alternatives,
  insatisfaction \\
\f{Conseil} & Conseiller & propositions (\f{Proposition} : produit, raisons, sources, score),
  écartés (\f{Ecart} : produit, raison, règle --- le \og pourquoi pas \fg), alertes, faits (les
  seuls nombres autorisés dans la réponse) \\
\f{Action} & Agir & outil, paramètres, confirmée (faux par défaut), résultat \\
\f{Etape} & Mesurer & une ligne du panneau Frugalité : nom, modèle, latence, tokens, appel au LLM \\
\f{Reponse} & \f{reply()} & texte, compréhension, conseil, action, question posée, escalade, trace \\
\f{EtatConversation} & \f{reply()} & la mémoire entre deux messages : parcours, information
  attendue, informations recueillies, action en attente, mots simples, historique \\
\bottomrule
\end{tabularx}
\end{center}

\subsection{Le chemin d'un message dans le code}

Tout part de \f{reply(message, client\_id, etat)}, dans \f{src/assistant.py} :

\begin{code}
\footnotesize
\begin{verbatim}
profil = clients.profil(client_id)                        # conseiller/clients.py -> clients.json

# 0. Une action attend un "oui" ?          -> outils.executer(action)   (conseiller/outils.py)

# 1. Comprendre                                                  (comprendre/__init__.py)
comprehension = comprendre(message, etat)
#    etat.info_attendue ?            -> c'est la reponse a la question posee
#    mots_cles.detecter(message)     -> sur ? on s'arrete
#    classifieur.predire(message)    -> au-dessus de tau ? on s'arrete
#    + entites.extraire(message)     -> montants, durees...
#    pas compris -> Decision("clarifier"), puis Decision("orienter") + incompris.noter()

# 2. Converser                                                   (converser/dialogue.py)
decision = avancer(message, comprehension, etat, profil)
#    lit parcours/<intention>.yaml, remplit etat.slots (profil, puis message)
#    -> "question" | "conseil" | "action" | "transmettre" | "fiche" | "abandon"

# 3. Conseiller (si conseil)                                     (conseiller/regles.py)
conseil = conseiller(etat.parcours, profil, etat.slots)

# 4. Agir (si action) : sans effet de bord -> execute ; sinon -> attend un "oui"

# 5. Formuler                                                    (converser/formuler.py)
texte = formuler(decision, conseil, action, profil, etat, message)
#    gabarit -> [LLM + garde-fou si ASSISTANT_LLM=1 et conseil] -> mots simples -> mention legale

return Reponse(texte, comprehension, conseil, action, ..., trace)   # chaque etape chronometree
\end{verbatim}
\end{code}

\subsection{Les briques une par une}

\paragraph{\f{mots\_cles.py} --- le filtre d'entrée.} Un dictionnaire
\f{\{intention: (expressions)\}}, cherché dans le message normalisé (\og VOLÉ MA CB \fg{}
devient \og vole ma cb \fg). Règle \og sûr \fg : une seule intention trouvée, et soit 12 mots au
plus, soit au moins 2 expressions.

\paragraph{\f{entites.py} --- les informations du message.} Des expressions régulières :
\og 250\,000 \fg, \og 250k \fg{} et \og 1\,900\eur{} \fg{} deviennent des montants ; \og 10 ans
\fg{} donne 120 mois ; \og mon fils \fg{} donne le lien ; \og oui \fg{} et \og non \fg{} pour les
confirmations. Durées et heures sont retirées avant de chercher les montants.

\paragraph{\f{parcours/*.yaml} et \f{dialogue.py} --- l'automate.} Un parcours, par exemple
\f{credit\_immobilier.yaml} (extrait) :

\begin{code}
\footnotesize
\begin{verbatim}
infos:
  revenus_nets:
    type: montant
    source: revenus_nets                     # lu dans le profil...
    confirmer: "J'ai note vos revenus de {valeur} nets par mois."   # ...et confirme
  prix_bien:
    type: montant
    indices: [prix, bien]                    # pour les corrections
    question: "Quel est le prix du bien que vous visez ?"
fin: {type: conseil, mention_legale: credit, proposer: transmettre_conseiller}
\end{verbatim}
\end{code}

\f{avancer()} remplit les informations depuis le profil (\f{source}) puis depuis le message
(selon le \f{type}), pose la première question qui manque, sinon déclenche la fin. Clés
spéciales : \f{si\_inconnu} (valeur prise sur \og je ne sais pas \fg), \f{si\_non} (valeur prise
sur \og non \fg), \f{indices} (à quelle information va un montant lors d'une correction).

\paragraph{\f{regles.py} --- le moteur de conseil.} Une fonction par parcours. Le crédit
immobilier applique la formule d'annuité, capacité $= m \times \frac{1-(1+i)^{-n}}{i}$, avec les
paramètres du catalogue (35\,\%, 25 ans, frais de notaire, taux indicatif). Il rend un
\f{Conseil} avec raisons, écarts et identifiant de règle (\f{HCSF-taux-effort}).

\paragraph{\f{outils.py} --- les actions.} \f{OUTILS} déclare chaque outil et s'il a un effet
de bord. \f{executer()} refuse un outil à effet de bord non confirmé (exception
\f{ConfirmationRequise}). Le SI est simulé sur une copie du profil ; chaque action est notée dans
un journal.

\paragraph{\f{formuler.py} et \f{garde\_fou.py} --- la réponse.} Un gabarit par type de décision
et par outil. Si \f{ASSISTANT\_LLM=1} et que c'est un conseil, un seul appel à Ollama reformule
le gabarit ; \f{verifier()} compare les nombres de la reformulation à ceux du gabarit et
refuse les promesses (\og accordé \fg, \og garanti \fg). Il refuse aussi les chiffres
\emph{oubliés} et les réponses tronquées : un test réel avec \f{gemma3:270m} répondait \og Oui,
je suis prêt à reformuler le message \fg{} sans le conseil. Au moindre écart, le gabarit est gardé.
La même fonction \f{reformuler()} sert à l'assistant et au banc de formulation.

\paragraph{\f{glossaire.py} --- le jargon.} Recherche exacte du terme, puis, pour une question
de définition, recherche approchante (\f{difflib}) qui rattrape les fautes : \og TEAG \fg{} donne
TAEG. Pas de modèle : une vingtaine de fiches écrites et relues.

\paragraph{\f{app.py} --- l'interface.} Barre latérale : choix du client (C01 à C08), messages de
démo cliquables, interrupteurs \og Mots simples \fg{} et \og Grand texte \fg. Trois zones : la
conversation, \og Pourquoi \fg{} (raisons, écarts, sources) et \og Frugalité \fg{} (une ligne par
étape, totaux de la conversation). La barre latérale affiche l'état \emph{réel} du LLM (serveur
interrogé, modèle préchargé), pas seulement la configuration.

\subsection{La configuration}

\begin{center}\small
\begin{tabularx}{\textwidth}{@{}L{4.5cm} L{4.1cm} Y@{}}
\toprule
\textbf{Variable} & \textbf{Défaut} & \textbf{Effet} \\
\midrule
\f{ASSISTANT\_LLM} & \f{0} & \f{1} : le petit LLM reformule les conseils \\
\f{ASSISTANT\_ESCALADE} & \f{0} & \f{1} : appel au 14B si un message reste incompris \\
\f{LLM\_API} (et \f{\_FORMULATION}, \f{\_ESCALADE}) & \f{ollama} & \f{ollama} ou \f{vllm}, par rôle \\
\f{VLLM\_BASE\_URL} & --- & Serveur vLLM de DeepEirb (\f{\dots/v1}) \\
\f{OLLAMA\_BASE\_URL} & \f{localhost:11434} & Adresse d'Ollama \\
\f{OLLAMA\_MODEL} & \f{ministral-3:3b} & Modèle de formulation \\
\f{OLLAMA\_MODEL\_ESCALADE} & \f{ministral-3:14b} & Modèle d'escalade \\
\f{INCOMPRIS\_FICHIER} & \f{eval/incompris.jsonl} & Où écrire le journal \\
\bottomrule
\end{tabularx}
\end{center}

\subsection{Lancer, tester, évaluer}

\begin{code}
\footnotesize
\begin{verbatim}
streamlit run app.py                               # l'interface
python -m pytest -q                                # 171 tests
python -m eval.run_scenarios                       # 9 conversations de bout en bout
python -m eval.bench_formulation --situations      # le banc LLM (sans modele : les 30 cas)
python -m src.incompris                            # relire les messages non compris
python -m eval.demo_secours                        # rejoue la demo -> rapport/demo-secours.tex
python -m eval.tableau_formulation                 # resultats du banc -> diapositive
\end{verbatim}
\end{code}

\textbf{Les tests} vérifient chaque brique : la taxonomie couvre les 77 intentions officielles ;
les mots-clés reconnaissent les situations des slides, et \og carte \fg{} seul ne suffit pas ; le
calcul HCSF donne 1\,197\eur{} et 41,3\,\% ; les actions refusent de s'exécuter sans
confirmation ; le garde-fou refuse les nombres inventés ; l'interface démarre et répond.

\textbf{Les scénarios} rejouent les conversations des 8 clients et vérifient ce qui est demandé,
ce qui ne doit \emph{pas} l'être, les produits proposés et écartés, et deux règles partout : au
plus un appel au LLM par tour, aucune action sans confirmation.

\end{document}
